\section{What migrant women describe}
\subsection{Domestic payments can work while a cross-border transfer still needs an intermediary}

Women described using wallets, bank transfers and QR payments to buy goods, receive customer payments and manage small businesses. In the accounts of sending money to Venezuela, a transfer from a Colombian wallet could start rather than finish the transaction. Women described asking trusted contacts or exchange operators for rates, transferring money locally, waiting for the other leg, and checking a receipt or asking relatives to confirm delivery. The legacy chapter also reports more direct or faster experiences in some other corridors; dependence on intermediaries was not uniform. These accounts concern different journeys and sometimes the opposite direction from the survey's inbound questions.

An intermediary could provide a service the local wallet did not: arrange the rate, coordinate delivery on the other side and confirm receipt. Some women used recommendations, repeat dealings and comparisons among contacts to decide whom to trust. These arrangements made transfers possible and sometimes offered useful reassurance. They also left women with work that did not end when the local transfer settled. Where charges and exchange-rate deductions were hard to separate, a fast local payment could still be part of a costly or time-consuming journey. These are reported experiences, not a price audit or a quantified loss of remittance value.

\subsection{Documents, practical help and control over the account}

Some women described trouble opening a product because they lacked a requested document; others said a document they held was not recognized by a provider or an app. Some relied on a relative's or partner's account. The same assistance could enable a transaction and limit privacy or direct control over it. Women described learning through use, comparing rates and helping others transact. The accounts tie these practices to livelihoods and to obligations such as paying for food, medicine or schooling across borders. Knowing how to move money and deciding where it should go do not guarantee control over the account or timely delivery.

\subsection{Taking precautions is not the same as having recourse}

The interviews and group discussions describe suspicious links, false confirmations and uncertainty about responsibility when a transfer was delayed or went wrong. Women reported checking requests and receipts, calling relatives and returning to known contacts. They could take these precautions and still find it hard to locate a complaint channel or know whether funds could be recovered. The legacy analysis distinguishes women facing barriers to formal entry, users who depend on intermediaries, digitally capable users constrained by available services, and women who transact more independently but still face corridor constraints. It treats security concerns as cross-cutting. These positions can overlap or change. They are not linked to survey respondents or assigned quantitative class labels.
